\documentclass[11pt]{article}

\usepackage[margin=1in]{geometry}
\usepackage{amsmath,amssymb,amsthm}

\newtheorem{theorem}{Theorem}
\newtheorem{lemma}{Lemma}
\newtheorem{proposition}{Proposition}

\newcommand{\Aut}{\operatorname{Aut}}
\newcommand{\dom}{\operatorname{dom}}

\begin{document}
\title{Submission C solution to Problem 1}
\date{}
\maketitle


\begin{abstract}
We answer the question affirmatively.  We build a computable countable
structure $\mathcal A$ such that every automorphism of every copy of
$\mathcal A$ is computable from that copy, but for every finite tuple of
parameters $\bar a$ there is an automorphism $\pi$ whose graph is not
$\Sigma^{\mathrm{in}}_1$-definable over $\bar a\cup\pi(\bar a)$.
\end{abstract}

\section*{The construction}

Let
\[
G=[\omega]^{<\omega}
\]
be the countable Boolean group of finite subsets of $\omega$, with operation
symmetric difference $\triangle$.  For $g\in G$ and $n\in\omega$, write
\[
\epsilon_n(g)=
\begin{cases}
1,& n\in g,\\
0,& n\notin g.
\end{cases}
\]
Let $e_n=\{n\}$.

We use a computable countable language $L$ containing:
\begin{itemize}
\item unary predicates $U,V_0,V_1,\dots$;
\item for each $g\in G$, a binary relation symbol $Q_g$;
\item for each $n\in\omega$, a binary relation symbol $P_n$;
\item one unary function symbol $S$.
\end{itemize}

The domain of $\mathcal A$ is the disjoint union
\[
U^{\mathcal A}\ \cup\ \bigcup_{n\in\omega}V_n^{\mathcal A},
\]
where
\[
U^{\mathcal A}=G
\]
and
\[
V_n^{\mathcal A}=V_n^0\cup V_n^1,\qquad
V_n^i=\{v_{n,i,m}:m\in\omega\}.
\]
The predicates $U,V_n$ name these sorts.  The symbols are interpreted as follows:
\[
Q_g(u,v)\quad\Longleftrightarrow\quad u,v\in U^{\mathcal A}
\ \text{ and }\ v=u\triangle g;
\]
\[
P_n(u,v_{n,i,m})\quad\Longleftrightarrow\quad
u\in U^{\mathcal A}\ \text{ and }\ i=\epsilon_n(u),
\]
and $P_n$ is false outside $U^{\mathcal A}\times V_n^{\mathcal A}$.  Finally,
\[
S(u)=u\quad(u\in U^{\mathcal A}),
\]
and
\[
S(v_{n,i,m})=v_{n,i,m+1}.
\]
With any standard computable coding of $G$ and of the triples $(n,i,m)$, this is a computable presentation of $\mathcal A$.

\section*{The automorphism group}

For each $g\in G$, define a permutation $\tau_g$ of $\mathcal A$ by
\[
\tau_g(u)=g\triangle u\qquad(u\in U^{\mathcal A}),
\]
and
\[
\tau_g(v_{n,i,m})=v_{n,i\oplus \epsilon_n(g),m}.
\]
Here $\oplus$ denotes addition modulo $2$.

\begin{lemma}
For every $g\in G$, $\tau_g\in\Aut(\mathcal A)$, and every automorphism of
$\mathcal A$ is of this form.  Hence
\[
\Aut(\mathcal A)=\{\tau_g:g\in G\}\cong G
\]
is countable.
\end{lemma}

\begin{proof}
It is immediate from the definitions that $\tau_g$ preserves all sort predicates and preserves $S$.

If $Q_h(u,v)$ holds, then $v=u\triangle h$, and therefore
\[
\tau_g(v)=g\triangle u\triangle h=\tau_g(u)\triangle h,
\]
so $Q_h(\tau_g(u),\tau_g(v))$ holds.

Similarly, if $P_n(u,v_{n,i,m})$ holds, then $i=\epsilon_n(u)$.  Since
\[
\epsilon_n(g\triangle u)=\epsilon_n(g)\oplus\epsilon_n(u),
\]
we have
\[
P_n(\tau_g(u),\tau_g(v_{n,i,m})).
\]
Thus $\tau_g$ is an automorphism.

Conversely, let $\alpha\in\Aut(\mathcal A)$.  Since $U$ is named, $\alpha$ maps $U^{\mathcal A}$ to itself.  Put
\[
g=\alpha(\varnothing)\in G.
\]
For every $u\in G$, the relation $Q_u(\varnothing,u)$ holds.  Applying $\alpha$, we get
\[
Q_u(g,\alpha(u)),
\]
so
\[
\alpha(u)=g\triangle u.
\]

Now fix $n$ and consider $V_n$.  Let $v=v_{n,i,m}$.  Choose $u\in G$ with
$\epsilon_n(u)=i$.  Then $P_n(u,v)$ holds.  Applying $\alpha$, we get
\[
P_n(g\triangle u,\alpha(v)).
\]
Therefore $\alpha(v)$ lies in the copy $V_n^{i\oplus\epsilon_n(g)}$.

Finally, $\alpha$ preserves $S$.  On each copy $V_n^i$, the map $S$ is the successor map on a one-way chain of type $\omega$.  Any bijection between two such chains preserving successor must send the unique element not in the range of $S$ to the unique element not in the range of $S$, and hence must preserve the coordinate $m$.  Thus
\[
\alpha(v_{n,i,m})=v_{n,i\oplus\epsilon_n(g),m}.
\]
So $\alpha=\tau_g$.
\end{proof}

\section*{Every automorphism of every copy is relatively computable}

\begin{proposition}
For every copy $\mathcal B\cong\mathcal A$ and every
$\beta\in\Aut(\mathcal B)$, the function $\beta$ is computable relative to
the atomic diagram of $\mathcal B$.  In particular, $\mathcal A$ is
computably $\Aut$-countable on the cone with base $\emptyset$.
\end{proposition}

\begin{proof}
Let $\mathcal B$ be any copy of $\mathcal A$, and let
$\beta\in\Aut(\mathcal B)$.  By the previous lemma, transferred to
$\mathcal B$, there is a unique $g\in G$ such that $\beta$ acts on the $U$-sort as translation by $g$.  Since $g$ is a finite subset of $\omega$, it can be hard-coded into a Turing program computing $\beta$.

Choose any $u_0\in U^{\mathcal B}$ and let
\[
u_1=\beta(u_0).
\]
The program may hard-code the natural numbers $u_0,u_1$.

Given input $x\in U^{\mathcal B}$, search through $h\in G$ until finding the unique $h$ such that
\[
\mathcal B\models Q_h(u_0,x).
\]
Then search for the unique $y$ such that
\[
\mathcal B\models Q_h(u_1,y).
\]
Output this $y$.  This computes $\beta(x)$ for $x\in U^{\mathcal B}$.

Now suppose $x\in V_n^{\mathcal B}$.  We can find $n$ by searching the unary predicates $V_0,V_1,\dots$.

If $n\notin g$, then $\beta$ is the identity on $V_n^{\mathcal B}$, so output $x$.

If $n\in g$, then $\beta$ swaps the two $S$-chains inside $V_n^{\mathcal B}$.  Choose some
\[
p_n\in V_n^{\mathcal B}
\]
and let
\[
q_n=\beta(p_n).
\]
The program may hard-code the finitely many pairs $(p_n,q_n)$ for $n\in g$.

To compute $\beta(x)$ for $x\in V_n^{\mathcal B}$ with $n\in g$, dovetail the following searches:
\begin{itemize}
\item search for $d$ with $S^d(p_n)=x$; if found, output $S^d(q_n)$;
\item search for $d$ with $S^d(x)=p_n$; if found, search for the unique $y$ with $S^d(y)=q_n$, and output $y$;
\item search for $d$ with $S^d(q_n)=x$; if found, output $S^d(p_n)$;
\item search for $d$ with $S^d(x)=q_n$; if found, search for the unique $y$ with $S^d(y)=p_n$, and output $y$.
\end{itemize}
Exactly one of these alternatives eventually applies, because each $V_n$ consists of two one-way $S$-chains and $p_n,q_n$ are corresponding points in the two chains.

All searches use only the atomic diagram of $\mathcal B$.  Thus $\beta$ is computable relative to $\mathcal B$.
\end{proof}

\section*{A preservation lemma for existential formulas}

Fix $n$.  For $d_0,d_1\in\omega$, define a map
\[
\rho^{n}_{d_0,d_1}:\mathcal A\to\mathcal A
\]
by fixing every element outside $V_n$, and by
\[
\rho^{n}_{d_0,d_1}(v_{n,i,m})=v_{n,i,m+d_i}.
\]

\begin{lemma}
For every $n$ and every $d_0,d_1\in\omega$, the map
$\rho^{n}_{d_0,d_1}$ is an embedding of $\mathcal A$ into itself.
Consequently, if an existential formula with parameters outside $V_n$
holds of a tuple, then it still holds after applying $\rho^{n}_{d_0,d_1}$
to the non-parameter entries.
\end{lemma}

\begin{proof}
The map $\rho^{n}_{d_0,d_1}$ is injective.  It clearly preserves the unary predicates.  It fixes $U$, so it preserves every $Q_g$.  It also preserves $S$, because
\[
\rho^{n}_{d_0,d_1}(S(v_{n,i,m}))
=
v_{n,i,m+1+d_i}
=
S(\rho^{n}_{d_0,d_1}(v_{n,i,m})).
\]
Finally, $P_n(u,v_{n,i,m})$ depends only on $u$ and on the copy index $i$, and $\rho^{n}_{d_0,d_1}$ fixes $u$ and preserves $i$.  For $m\neq n$, $P_m$ is unaffected.

Thus $\rho^{n}_{d_0,d_1}$ is an $L$-embedding.  Embeddings preserve quantifier-free formulas, and therefore preserve existential formulas.
\end{proof}

\section*{Failure of $\Sigma^{\mathrm{in}}_1$-definability}

\begin{lemma}
Let $\bar p$ be a finite tuple from $\mathcal A$ containing no element of
$V_n$.  Then the automorphism $\tau_{e_n}$ is not
$\Sigma^{\mathrm{in}}_1$-definable over $\bar p$.
\end{lemma}

\begin{proof}
Let
\[
r_0=v_{n,0,0},\qquad r_1=v_{n,1,0}.
\]
Then
\[
\tau_{e_n}(r_0)=r_1.
\]

Suppose, toward a contradiction, that the graph of $\tau_{e_n}$ is
$\Sigma^{\mathrm{in}}_1$-definable over $\bar p$.  Then there is a countable
family of existential formulas
\[
(\varphi_i(\bar p,x,y))_{i\in\omega}
\]
such that
\[
\tau_{e_n}(x)=y
\quad\Longleftrightarrow\quad
\mathcal A\models \varphi_i(\bar p,x,y)
\ \text{ for some }i.
\]
Since $\tau_{e_n}(r_0)=r_1$, there is some $i$ with
\[
\mathcal A\models \varphi_i(\bar p,r_0,r_1).
\]

Now apply the embedding $\rho^n_{0,1}$.  It fixes $\bar p$ and $r_0$, and sends
\[
r_1\mapsto S(r_1)=v_{n,1,1}.
\]
By preservation of existential formulas under embeddings,
\[
\mathcal A\models \varphi_i(\bar p,r_0,S(r_1)).
\]
But
\[
\tau_{e_n}(r_0)=r_1\neq S(r_1),
\]
so $(r_0,S(r_1))$ is not in the graph of $\tau_{e_n}$.  This contradicts the assumed definition of the graph.
\end{proof}

\section*{Main theorem}

\begin{theorem}
There is a countable computably represented structure $\mathcal A$ such that:
\begin{enumerate}
\item $\mathcal A$ is computably $\Aut$-countable on a cone;
\item for every finite tuple of parameters $\bar a$ from $\mathcal A$, there is an automorphism $\pi\in\Aut(\mathcal A)$ such that $\pi$ is not
$\Sigma^{\mathrm{in}}_1$-definable over $\bar a\cup\pi(\bar a)$.
\end{enumerate}
\end{theorem}

\begin{proof}
The structure $\mathcal A$ constructed above is computable and countable.  By the proposition, every automorphism of every copy of $\mathcal A$ is computable relative to that copy, so $\mathcal A$ is computably $\Aut$-countable on the cone with base $\emptyset$.

Now let $\bar a$ be any finite tuple from $\mathcal A$.  Since $\bar a$ is finite, it meets only finitely many of the sorts $V_n$.  Choose $n$ such that no entry of $\bar a$ lies in $V_n$.

Let
\[
\pi=\tau_{e_n}.
\]
Then $\pi$ toggles the two copies inside $V_n$ and fixes every $V_m$ with $m\neq n$.  Since $\bar a$ has no entry in $V_n$, the tuple $\pi(\bar a)$ also has no entry in $V_n$.  Therefore the finite tuple
\[
\bar p=\bar a\cup\pi(\bar a)
\]
contains no element of $V_n$.

By the previous lemma, the graph of $\pi=\tau_{e_n}$ is not
$\Sigma^{\mathrm{in}}_1$-definable over $\bar p$.  Hence $\pi$ is not
$\Sigma^{\mathrm{in}}_1$-definable over $\bar a\cup\pi(\bar a)$.

Thus $\mathcal A$ has the required properties.
\end{proof}

\end{document}